Write the outside face residual as
\begin{equation*}
 \vr^\gU:=\vc_\rho-\mQ\vx^\gU=-\vw^\gU.
\end{equation*}
Fix $\vartheta\geq0$.  A nested sequence
$\gU_k=\gB_0\mathbin{\dot\cup}\cdots\mathbin{\dot\cup}\gB_k$ is
\emph{$\vartheta$-batched} if $\gB_0=\{v\}$, every later $\gB_k$ is nonempty
and satisfies $r_i^{\gU_{k-1}}>0$ for every $i\in\gB_k$, and, immediately
before admitting $\gB_k$,
\[
 r_i^{\gU_{k-1}}\leq\vartheta
 \qquad\text{for every }i\in\gS_\rho^*\setminus\gU_k.
\]
The condition includes coordinates not yet assigned to a later block.
Any finite such sequence can be extended to partition $\gS_\rho^*$ by
repeatedly admitting all exact positive residuals. After each added batch,
every remaining coordinate had nonpositive residual on the preceding face,
so the threshold condition is preserved. This is an analytical condition;
the algorithm need not know $\gS_\rho^*$. The residual certificate will
establish it for all unreported coordinates in
\cref{sec:threshold-batch-analysis}.

\begin{lemma}[Chebyshev inverse approximation]
\label{lem:chebyshev-inverse}
Let $\alpha\mI\preceq\mM\preceq b_0\mI$, with $b_0\geq\alpha>0$.
For each $j\geq1$, there is a polynomial $p_{j-1}$ of degree at most
$j-1$ such that
\begin{equation*}
 \norm{\mM^{-1}-p_{j-1}(\mM)}_2
 \leq\frac2\alpha
       \left(\frac{\sqrt{b_0/\alpha}-1}
                  {\sqrt{b_0/\alpha}+1}\right)^j.
\end{equation*}
The corresponding residual polynomial has absolute value at most twice
the displayed geometric factor on $[\alpha,b_0]$.
\end{lemma}
\begin{proof}
If $b_0=\alpha$, take $p_{j-1}(t)=1/\alpha$; both errors vanish.
Otherwise $b_0>\alpha$.
Let $T_j$ be the first-kind Chebyshev polynomial and set
\[
 r_j(t)=
 \frac{T_j((b_0+\alpha-2t)/(b_0-\alpha))}
      {T_j((b_0+\alpha)/(b_0-\alpha))}.
\]
Its numerator has absolute value at most one on the interval.
For $u>1$,
$T_j(u)=((u+\sqrt{u^2-1})^j+(u-\sqrt{u^2-1})^j)/2$.
Substitution bounds $|r_j(t)|$ by twice the stated geometric factor.
Also $r_j(0)=1$, so $p_{j-1}(t)=(1-r_j(t))/t$ is a polynomial of
degree at most $j-1$ and
$|1/t-p_{j-1}(t)|=|r_j(t)|/t\leq|r_j(t)|/\alpha$.
Diagonalizing the symmetric matrix proves the operator-norm bound.
\end{proof}

\begin{theorem}[Threshold-batch energy-depth bound]
\label{thm:threshold-batch-depth}
Let $\gU_0\subset \gU_1\subset\cdots$ be a $\vartheta$-batched reachable
sequence, extended to the full optimal support.  Put
\begin{equation*}
 q_\alpha:=
 \frac{\sqrt{2/\alpha}-1}{\sqrt{2/\alpha}+1}.
\end{equation*}
Then every included face $\gU_J$ satisfies
\begin{equation}
 \phi_\rho(\vx^{\gU_J})-\phi_\rho(\vx_\rho^*)
 \leq8q_\alpha^{2J}+\frac{\vartheta^2}{\alpha\rho}.
 \label{eq:batch-depth-main}
\end{equation}
For $0<\epsobj<1$, the decaying term is at most $\epsobj/2$ after
$O(\alpha^{-1/2}\log(2/\epsobj))$ batches. If the sequence ends sooner,
its final face solution is exact. The estimate is independent of
$|\gS_\rho^*|$, graph diameter, and any complementarity margin.
\end{theorem}

The proof has three steps. Block Cholesky elimination expresses the
objective gap as a tail energy. A block-tridiagonal comparison system
preserves spectral bounds and encodes delayed admissions. Chebyshev
approximation then separates the decaying seed contribution from the
threshold error. These factors are analytical objects and are not
constructed by the local algorithm.

\begin{proof}
Let $\gS=\gS_\rho^*$ and order it by the blocks
$\gB_0,\gB_1,\ldots,\gB_K$.  Block elimination gives
\begin{equation*}
 \mQ_{\gS\gS}=\mL\mDelta\mL^\trans,
 \qquad
 \mDelta=\operatorname{diag}(\mDelta_0,\ldots,\mDelta_K),
 \qquad
 \vg=\mL^{-1}(\vc_\rho)_\gS.
\end{equation*}
Here $\mL$ is unit block lower triangular, every pivot $\mDelta_k$ is
Stieltjes, and $\vg_k$ is the positive Schur residual on $\gB_k$ when it is
admitted. Let $\mR_k$ be its lower Cholesky factor, so
$\mDelta_k=\mR_k\mR_k^\trans$, and set
\begin{equation*}
 \vz_k:=\mR_k^{-1}\vg_k.
\end{equation*}
Each $\mR_k$ is a triangular nonsingular M-matrix, so $\vz_k\geq\vzero$.
Completion of squares in the block elimination gives the exact identity
\begin{equation}
 2\bigl(\phi_\rho(\vx^{\gU_J})-\phi_\rho(\vx_\rho^*)\bigr)
 =\sum_{k>J}\norm{\vz_k}_2^2.
 \label{eq:batch-exact-gap-tail}
\end{equation}

Put
$\mC=\mL\operatorname{diag}(\mR_0,\ldots,\mR_K)$, so that
$\mQ_{\gS\gS}=\mC\mC^\trans$.  This is the scalar lower Cholesky factor in the
block-respecting order.  Its recursion shows that every off-diagonal entry is
nonpositive.  Let $\widehat{\mC}$ retain only the diagonal blocks and first
block subdiagonal.  Both factors are triangular nonsingular M-matrices and
$\mC\leq\widehat{\mC}$ entrywise.  Hence
\begin{equation}
 \vzero\leq\widehat{\mC}^{-1}\leq\mC^{-1},
 \qquad
 \norm{\widehat{\mC}^{-1}}_2
 \leq\norm{\mC^{-1}}_2\leq\frac1{\sqrt\alpha}.
 \label{eq:batch-bidiagonal-inverse}
\end{equation}
The inverse ordering follows directly from
\[
 \mC^{-1}-\widehat{\mC}^{-1}
 =\mC^{-1}(\widehat{\mC}-\mC)\widehat{\mC}^{-1}\geq\vzero.
\]
The spectral-norm comparison is valid because both inverses are nonnegative:
$|\widehat{\mC}^{-1}\vz|\leq\mC^{-1}|\vz|$ for every $\vz$.

There is also a graph-independent upper bound.  The $k$th block row of
$\mC$ satisfies
\begin{equation*}
 \sum_{i\leq k}\mC_{ki}\mC_{ki}^\trans
 =\mQ_{\gB_k\gB_k}\preceq\mI.
\end{equation*}
Deleting all but two input blocks cannot increase its row-map norm, and each
input block occurs in at most two retained rows.  Therefore
\begin{equation*}
 \norm{\widehat{\mC}}_2\leq\sqrt2.
\end{equation*}
Consequently
$\mH_{\mathrm{ch}}:=\widehat{\mC}\widehat{\mC}^\trans$ is block tridiagonal and
\begin{equation*}
 \alpha\mI\preceq\mH_{\mathrm{ch}}\preceq2\mI.
\end{equation*}

We next encode causality.  For $k\geq2$, immediately before $\gB_{k-1}$ is
admitted, the still-future block $\gB_k$ has residual at most
$\vartheta\one$; this is its residual after elimination through $\gB_{k-2}$.
For $k=1$, the same bound follows from
$(\vc_\rho)_i=-\alpha\rho\sqrt{d_i}<0$ away from the seed.  Eliminating
$\gB_{k-1}$ subtracts $\mL_{k,k-1}\vg_{k-1}$, so
\begin{equation*}
 \vg_k+\mL_{k,k-1}\vg_{k-1}\leq\vartheta\one
 \qquad(k\geq1).
\end{equation*}
Using
$\mC_{k,k-1}=\mL_{k,k-1}\mR_{k-1}$, this becomes
\begin{equation*}
 (\widehat{\mC}\vz)_0=\vg_0,
 \qquad
 (\widehat{\mC}\vz)_k\leq\vartheta\one\quad(k\geq1).
\end{equation*}
Since $\widehat{\mC}^{-1}\geq\vzero$,
\begin{equation}
 \vzero\leq\vz\leq
 \widehat{\mC}^{-1}
 \begin{bmatrix}\vg_0\\ \vartheta\one_{\gS\setminus \gB_0}\end{bmatrix}.
 \label{eq:batch-chain-domination}
\end{equation}

The distributed threshold source is bounded by
\begin{equation}
 \norm{\widehat{\mC}^{-1}[\vzero;\vartheta\one]}_2^2
 \leq\frac{\vartheta^2|\gS|}{\alpha}
 \leq\frac{\vartheta^2}{\alpha\rho},
 \label{eq:batch-threshold-source}
\end{equation}
where the last step uses
$|\gS|\leq\vol(\gS)\leq1/\rho$.

For the seed source and $J\geq1$, use \cref{lem:chebyshev-inverse} to take a
degree-$(J-1)$ polynomial
$p_{J-1}$ for $1/t$ on $[\alpha,2]$, with
\begin{equation*}
 \norm{\mH_{\mathrm{ch}}^{-1}-p_{J-1}(\mH_{\mathrm{ch}})}_2
 \leq\frac2\alpha q_\alpha^J.
\end{equation*}
Use $\widehat{\mC}^{-1}=\widehat{\mC}^{\trans}\mH_{\mathrm{ch}}^{-1}$.
Because $\mH_{\mathrm{ch}}$ is block tridiagonal,
$p_{J-1}(\mH_{\mathrm{ch}})[\vg_0;\vzero]$ reaches only blocks through
$J-1$. Multiplication by the block upper-bidiagonal
$\widehat{\mC}^\trans$ does not increase the largest block index.  If
$\mPi_{>J}$ projects onto the block tail, then
\begin{align}
 \norm{\mPi_{>J}\widehat{\mC}^{-1}[\vg_0;\vzero]}_2
 &\leq\norm{\widehat{\mC}}_2
       \norm{\mH_{\mathrm{ch}}^{-1}-p_{J-1}(\mH_{\mathrm{ch}})}_2
       \norm{\vg_0}_2\notag\\
 &\leq2\sqrt2q_\alpha^J.
 \label{eq:batch-seed-tail}
\end{align}
Here
$\vg_0=(\vc_\rho)_v=\alpha(1-\rho d_v)/\sqrt{d_v}\leq\alpha$.
For $J=0$, the same bound follows from
\cref{eq:batch-bidiagonal-inverse} and $\alpha\leq1$.
Combining
\cref{eq:batch-exact-gap-tail,eq:batch-chain-domination,eq:batch-threshold-source,eq:batch-seed-tail}
with $\norm{\va+\vb}_2^2/2\leq\norm{\va}_2^2+\norm{\vb}_2^2$ proves
\cref{eq:batch-depth-main}.
\end{proof}

The two terms have different meanings.  The seed term decays spectrally down
the block chain.  The threshold term charges every delayed coordinate once,
at the block where it is eventually eliminated.  A vertex can therefore
remain below threshold for many faces without generating a repeated,
uncharged error.
